(1) 质量恒定为 \(m\) 的理想完全反射光帆在一维真空中受从左向右照射的准直连续光束。固定实验室截面的入射束功率 \(P_0>0\)，光帆速度 \(v=\beta c\ge0\)。单个反射光子的镜面频率比给定为 \(\omega_r/\omega_i=(1-\beta)/(1+\beta)\)。特别注意移动光帆单位实验室时间截获光的功率并非 \(P_0\)。推导实际动量转移力 \(F(v)\)，给出静止光帆极限。
(2) 设初始静止 \(x(0)=0\)，忽略横向漏光、光帆吸收、束功率随距离变化和质量损失。引入快度 \(\chi=\operatorname{artanh}\beta\)，由 \(d(\gamma mv)/dt=F\) 求 \(t(\chi)\)、\(x(\chi)\) 的精确参数解及初始非相对论加速度。
(3) 求当 \(t\to\infty\) 时 \(\chi,\gamma,1-\beta\) 的时间渐近标度，并验证距离在任意有限时间内小于 \(ct\)，讨论理想无限准直光束的物理限制。
(4) 根据单位时间截获的入射能量和反射能量独立检验 \(d(\gamma mc^2)/dt=Fv\)，给出束能中进入机械动能的比例及接近光速时的行为；区别单个光子的反冲与准连续经典光强近似。
